\scanpage{52}
\mspage{II}
% archive foliation: 51
\uline{Lemme 2} Soit $X\to Y$ un morphisme de type fini, avec $Y$ loc. noeth. et univ. japonais.
Alors pour tt faisceau d'ensembles constructible $F$ sur $X$, $f_*(F)$ est constructible.

\unsure{§ A au début}

\uline{Lemme 3} $X$ schéma connexe \deleted{\unsure{loc}} noeth.,
\inserted{univ. japonais},
$(X_i)_{i\in I}$ une partition finie en parties loc. fermées, $G$ un groupe fini.
\deleted{Supposons que pour tout $i$, $H^1(X_i,G)$ soit fini.}
\inserted{Considérons l'\deleted{\unsure{application}}}
\deleted{Alors \unsure{pour tt}}
\diagram{\begin{tikzcd}
H^1(X,G)\arrow[r]&\displaystyle\prod_i H^1(X_i,G)
\end{tikzcd}}
Cette application est à \uline{fibres finies}. En particulier, si pour tt $i$, $H^1(X_i,G)$ est fini, kif kif pour $H^1(X,G)$.

\unsure{Théorie du recollement} des faisceaux étales, en utilisant lemme 2 et le fait suivant :
\inserted{\uline{Lemme 4}}
Si $F,G$ sont des faisceaux \inserted{d'ens} étales constr. sur $X$ noethérien, alors $\operatorname{Hom}(F,G)$ est un ens fini.
\margin{\figure{\begin{tikzpicture}
\draw (0,0)--(0,-1.2);
\end{tikzpicture}}}

\deleted{\unsure{Corollaire}}
\deleted{\unsure{Proposition}}
\inserted{\uline{Proposition}}
$S$ \deleted{\unsure{X}} schéma noeth. univ. japonais,
$G$ un groupe fini. Conditions équivalentes.
